\chapter{Design Patterns}
\chapterepigraph{A design pattern is a named, reusable answer to a
recurring design problem. You will almost never invent one in an
interview; you will \emph{recognise} that the problem in front of you is
one you have seen before and reach for the matching pattern. This chapter
is organised to make that recognition fast: learn the intent and the
tell-tale phrase of each, and the C++ writes itself.}

\section{How to Choose a Pattern}
\label{sec:patterns-intro}

The Gang of Four catalogue splits into three families by \emph{what the
pattern is about}: creating objects, composing objects into larger
structures, or coordinating behaviour between objects.

\begin{patternbox}[The Three Families and Their Job]
\begin{itemize}
  \item \textbf{Creational} -- control \emph{how objects are made}, so
        code does not hard-wire concrete classes: Singleton, Factory
        Method, Abstract Factory, Builder, Prototype.
  \item \textbf{Structural} -- compose objects and classes into larger
        structures while keeping them flexible: Adapter, Bridge,
        Composite, Decorator, Facade, Flyweight, Proxy.
  \item \textbf{Behavioural} -- assign responsibilities and manage
        communication between objects: Chain of Responsibility, Command,
        Iterator, Mediator, Memento, Observer, State, Strategy, Template
        Method, Visitor.
\end{itemize}
\end{patternbox}

\begin{ideabox}[Recognise by the Sentence the Interviewer Says]
Patterns are triggered by phrases. ``Exactly one instance'' $\to$
Singleton. ``Create objects without naming the class'' $\to$ Factory.
``Build a complex object step by step'' $\to$ Builder. ``Add behaviour at
run time without subclassing'' $\to$ Decorator. ``Make two incompatible
interfaces work together'' $\to$ Adapter. ``Notify many objects when one
changes'' $\to$ Observer. ``Swap an algorithm at run time'' $\to$
Strategy. ``Behaviour depends on internal state'' $\to$ State. Train on
the triggers and pattern choice becomes reflexive.
\end{ideabox}

Each section states the pattern's intent, the phrase that signals it, its
participants and structure, a complete C++ example, and the trade-offs
that decide whether it is worth the extra indirection.

\newpage

% ---- Creational ----
\section{Singleton}
\label{sec:pat-singleton}

\problemstatement{Ensure a class has exactly one instance and provide a
single global point of access to it -- for example a configuration store,
a logger, or a connection pool that must be shared, not duplicated.}

\begin{patternbox}[When You Reach for It]
The trigger phrase is \emph{``there must be exactly one''} -- one logger,
one config, one thread pool. Beware: Singleton is the most overused
pattern and effectively a controlled global, so use it only when a second
instance would be genuinely wrong, not merely inconvenient.
\end{patternbox}

\subsection*{Understanding the pattern}

Three moves enforce singularity: make the constructor \texttt{private} so
nobody can \texttt{new} it, delete the copy constructor and assignment so
it cannot be duplicated, and expose a static accessor that lazily creates
and returns the one instance. In C++ the cleanest, thread-safe form uses a
function-local \texttt{static}, whose initialisation the standard
guarantees to happen exactly once even under concurrency.

\begin{ideabox}[Meyers' Singleton: a Function-Local Static]
A \texttt{static} local variable inside the accessor is constructed on
first call and destroyed at program exit, and C++11 guarantees this
initialisation is thread-safe. This one-liner sidesteps the double-checked
locking that plagues Singleton in other languages.
\end{ideabox}

\subsection*{C++ implementation}

\begin{lstlisting}
#include <bits/stdc++.h>
using namespace std;

class Logger {
    Logger() { cout << "Logger created\n"; }     // private ctor
public:
    Logger(const Logger&) = delete;               // no copying
    Logger& operator=(const Logger&) = delete;

    static Logger& instance() {                   // single access point
        static Logger inst;                       // created once, thread-safe
        return inst;
    }
    void log(const string& msg) { cout << "[log] " << msg << "\n"; }
};

int main() {
    Logger::instance().log("first");
    Logger::instance().log("second");             // same instance
    cout << "&a == &b? "
         << (&Logger::instance() == &Logger::instance()) << "\n";
    return 0;
}
\end{lstlisting}

\begin{complexitybox}[Trade-offs]
\textbf{Pros.} Guarantees one instance; lazy creation; global access.
\textbf{Cons.} Introduces global state, which hides dependencies and makes
unit testing hard (you cannot easily substitute a mock). It can also
become a bottleneck or a hidden coupling point. Prefer dependency
injection of a single shared instance when testability matters.
\end{complexitybox}

\begin{takeawaybox}
Singleton = private constructor + deleted copies + a static accessor.
In modern C++ the function-local static gives you thread-safe lazy
initialisation for free. Reach for it sparingly; ``one instance'' is a
real constraint, but a shared injected object often serves the same need
without the testing pain of a global.
\end{takeawaybox}

\section{Factory Method}
\label{sec:pat-factory-method}

\problemstatement{Define an interface for creating an object, but let the
decision of \emph{which} concrete class to instantiate be made elsewhere,
so client code depends only on the abstract product, never on concrete
class names.}

\begin{patternbox}[When You Reach for It]
The trigger is \emph{``create an object without hard-coding its class,''}
usually to kill a \texttt{switch}/\texttt{if-else} on a type string that
\texttt{new}s different concrete classes. Whenever adding a new product
type forces you to edit existing creation code, a factory restores the
Open/Closed Principle.
\end{patternbox}

\subsection*{Understanding the pattern}

Client code should not contain \texttt{new ConcreteA} or \texttt{new
ConcreteB}. Instead it calls a factory function that takes a request
(often an enum or string) and returns a pointer to the abstract product.
The concrete decision lives in one place; the rest of the system speaks
only the interface. Adding a product means extending the factory, not
hunting down every \texttt{new} scattered across the codebase.

\begin{ideabox}[One Place Knows the Concrete Classes]
Centralise object creation behind a function returning the base type. The
caller says \emph{what} it wants (``a PDF exporter''); the factory decides
\emph{how} to build it. This single indirection is what lets you swap or
add implementations without touching callers.
\end{ideabox}

\subsection*{C++ implementation}

\begin{lstlisting}
#include <bits/stdc++.h>
using namespace std;

struct Shape {
    virtual void draw() const = 0;
    virtual ~Shape() = default;
};
struct Circle : Shape { void draw() const override { cout << "Circle\n"; } };
struct Square : Shape { void draw() const override { cout << "Square\n"; } };

enum class ShapeType { Circle, Square };

struct ShapeFactory {                          // the one place that knows
    static unique_ptr<Shape> create(ShapeType t) {
        switch (t) {
            case ShapeType::Circle: return make_unique<Circle>();
            case ShapeType::Square: return make_unique<Square>();
        }
        return nullptr;
    }
};

int main() {
    auto s1 = ShapeFactory::create(ShapeType::Circle);
    auto s2 = ShapeFactory::create(ShapeType::Square);
    s1->draw();                                // caller knows only "Shape"
    s2->draw();
    return 0;
}
\end{lstlisting}

\begin{complexitybox}[Trade-offs]
\textbf{Pros.} Decouples callers from concrete classes; new products are
added in one spot; supports the Open/Closed Principle. \textbf{Cons.} Adds
a layer of indirection and one more class/function to maintain; the
factory itself still centralises knowledge of every product (mitigated by
registration-based factories).
\end{complexitybox}

\begin{takeawaybox}
Factory Method replaces scattered \texttt{new ConcreteClass} calls with a
single creator returning the abstract type. It is the standard cure for a
type-tag \texttt{switch} and the most common way an interviewer expects
you to satisfy Open/Closed. Abstract Factory (next) generalises it to
whole \emph{families} of products.
\end{takeawaybox}

\section{Abstract Factory}
\label{sec:pat-abstract-factory}

\problemstatement{Provide an interface for creating \emph{families} of
related objects without specifying their concrete classes, so an entire
family can be swapped in one move -- for example a whole ``dark theme'' set
of UI widgets versus a ``light theme'' set.}

\begin{patternbox}[When You Reach for It]
The trigger is \emph{``several related objects that must be used
together and vary as a set.''} A Windows button must pair with a Windows
checkbox, not a Mac one. When products come in matched families and you
must guarantee consistency within a family, reach for Abstract Factory.
\end{patternbox}

\subsection*{Understanding the pattern}

Where Factory Method creates \emph{one} product, Abstract Factory creates a
\emph{suite}. You define an abstract factory interface with a creator
method per product type (\texttt{createButton}, \texttt{createCheckbox}).
Each concrete factory (\texttt{WindowsFactory}, \texttt{MacFactory})
produces the matching family. Client code holds an abstract factory and
asks it for products, never knowing or mixing platforms.

\begin{ideabox}[A Factory of Factories, One per Family]
Group the creators for a coherent family behind a single factory
interface. Selecting a concrete factory once (at start-up) fixes the whole
family, structurally preventing a mismatch like a Mac button beside a
Windows checkbox.
\end{ideabox}

\subsection*{C++ implementation}

\begin{lstlisting}
#include <bits/stdc++.h>
using namespace std;

struct Button   { virtual void paint() = 0; virtual ~Button()   = default; };
struct Checkbox { virtual void paint() = 0; virtual ~Checkbox() = default; };

struct WinButton   : Button   { void paint() override { cout << "Win button\n"; } };
struct WinCheckbox : Checkbox { void paint() override { cout << "Win checkbox\n"; } };
struct MacButton   : Button   { void paint() override { cout << "Mac button\n"; } };
struct MacCheckbox : Checkbox { void paint() override { cout << "Mac checkbox\n"; } };

struct GUIFactory {                            // abstract factory
    virtual unique_ptr<Button>   createButton()   = 0;
    virtual unique_ptr<Checkbox> createCheckbox() = 0;
    virtual ~GUIFactory() = default;
};
struct WinFactory : GUIFactory {
    unique_ptr<Button>   createButton()   override { return make_unique<WinButton>(); }
    unique_ptr<Checkbox> createCheckbox() override { return make_unique<WinCheckbox>(); }
};
struct MacFactory : GUIFactory {
    unique_ptr<Button>   createButton()   override { return make_unique<MacButton>(); }
    unique_ptr<Checkbox> createCheckbox() override { return make_unique<MacCheckbox>(); }
};

void renderUI(GUIFactory& f) {                 // platform-agnostic
    f.createButton()->paint();
    f.createCheckbox()->paint();
}

int main() {
    WinFactory win; MacFactory mac;
    renderUI(win);
    renderUI(mac);                             // whole family swapped
    return 0;
}
\end{lstlisting}

\begin{complexitybox}[Trade-offs]
\textbf{Pros.} Guarantees products from one family are used together;
swapping families is a one-line change; isolates concrete classes.
\textbf{Cons.} Adding a \emph{new product type} means changing the factory
interface and every concrete factory -- a real cost when the product set
grows. Best when families change often but the set of product types is
stable.
\end{complexitybox}

\begin{takeawaybox}
Abstract Factory is Factory Method scaled to families: one factory
interface with several creators, one concrete factory per family. It buys
family consistency at the price of rigidity when new product \emph{types}
appear. Recognise it whenever ``these objects must all match'' is a
requirement.
\end{takeawaybox}

\section{Builder}
\label{sec:pat-builder}

\problemstatement{Separate the construction of a complex object from its
representation, so the same step-by-step process can build different
configurations -- especially useful when an object has many optional
parameters.}

\begin{patternbox}[When You Reach for It]
The triggers are \emph{``an object with many optional fields''} (the
telescoping-constructor smell: \texttt{Pizza(true,false,true,...)}) and
\emph{``build a complex object in steps.''} When a constructor has a long
list of parameters, half of them optional, a Builder makes construction
readable and safe.
\end{patternbox}

\subsection*{Understanding the pattern}

Instead of one giant constructor, a Builder exposes chainable setter
methods (\texttt{addCheese()}, \texttt{addOlives()}) that accumulate
configuration, then a \texttt{build()} that produces the finished,
immutable product. Each setter returns the builder itself so calls read
like a sentence. The product's own constructor can stay private, taking
only the builder, guaranteeing objects are never half-initialised.

\begin{ideabox}[Chainable Steps, One Final Build]
Name each optional step explicitly and let the terminal \texttt{build()}
assemble the product. This turns an unreadable positional-argument call
into self-documenting, order-independent construction, and lets you
validate the whole configuration in one place before creating the object.
\end{ideabox}

\subsection*{C++ implementation}

\begin{lstlisting}
#include <bits/stdc++.h>
using namespace std;

class Pizza {
public:
    string size;
    bool cheese = false, olives = false, mushrooms = false;
    void describe() const {
        cout << size << " pizza"
             << (cheese ? " +cheese" : "")
             << (olives ? " +olives" : "")
             << (mushrooms ? " +mushrooms" : "") << "\n";
    }
};

class PizzaBuilder {
    Pizza p;
public:
    PizzaBuilder& setSize(const string& s) { p.size = s; return *this; }
    PizzaBuilder& addCheese()    { p.cheese = true;    return *this; }
    PizzaBuilder& addOlives()    { p.olives = true;    return *this; }
    PizzaBuilder& addMushrooms() { p.mushrooms = true; return *this; }
    Pizza build() { return p; }               // finished product
};

int main() {
    Pizza pizza = PizzaBuilder()
                    .setSize("Large")
                    .addCheese()
                    .addMushrooms()
                    .build();                 // reads like a sentence
    pizza.describe();
    return 0;
}
\end{lstlisting}

\begin{complexitybox}[Trade-offs]
\textbf{Pros.} Readable, order-independent construction; no telescoping
constructors; can enforce validation and immutability at \texttt{build}
time. \textbf{Cons.} More code than a plain constructor; overkill for
objects with a couple of fields. Worth it precisely when the parameter
list is long and largely optional.
\end{complexitybox}

\begin{takeawaybox}
Builder trades a long, error-prone constructor for a fluent chain of named
steps ending in \texttt{build()}. Recognise it by the smell of many
optional parameters or boolean flags at a call site. It is one of the most
frequently \emph{used} patterns in real code, not just interviews.
\end{takeawaybox}

\section{Prototype}
\label{sec:pat-prototype}

\problemstatement{Create new objects by cloning an existing prototype
rather than constructing from scratch, useful when object creation is
expensive or when the exact concrete type to copy is known only at run
time.}

\begin{patternbox}[When You Reach for It]
The trigger is \emph{``make a copy of an existing configured object,''} or
\emph{``creating one from scratch is costly.''} When you have a
fully-configured instance and want more just like it -- without knowing or
naming its concrete class -- a polymorphic \texttt{clone()} is the answer.
\end{patternbox}

\subsection*{Understanding the pattern}

Each class implements a virtual \texttt{clone()} that returns a copy of
itself as the base type. Client code holds a base pointer and calls
\texttt{clone()}; it gets back a new object of the correct concrete type
without a \texttt{switch} on types or a call to a specific constructor.
This is invaluable when the set of concrete classes is registered at run
time (a palette of prototype shapes the user can duplicate).

\begin{ideabox}[Let the Object Copy Itself]
Because only the object knows its own concrete type, delegate copying to a
virtual \texttt{clone()}. This sidesteps ``which constructor do I call?''
and preserves all the run-time state of the prototype, including fields a
factory would not know to set.
\end{ideabox}

\subsection*{C++ implementation}

\begin{lstlisting}
#include <bits/stdc++.h>
using namespace std;

struct Shape {
    int x = 0, y = 0;
    virtual unique_ptr<Shape> clone() const = 0;   // polymorphic copy
    virtual void draw() const = 0;
    virtual ~Shape() = default;
};

struct Circle : Shape {
    int r = 1;
    unique_ptr<Shape> clone() const override {
        return make_unique<Circle>(*this);          // copy all state
    }
    void draw() const override {
        cout << "Circle r=" << r << " at (" << x << "," << y << ")\n";
    }
};

int main() {
    Circle base; base.r = 5; base.x = 2; base.y = 3;

    unique_ptr<Shape> proto = make_unique<Circle>(base);
    auto copy1 = proto->clone();                    // no concrete type named
    auto copy2 = proto->clone();
    copy1->draw();
    copy2->draw();
    return 0;
}
\end{lstlisting}

\begin{complexitybox}[Trade-offs]
\textbf{Pros.} Copies run-time state exactly; avoids costly
re-construction; adds objects without naming concrete classes.
\textbf{Cons.} Deep vs.\ shallow copy must be handled carefully
(pointers, cycles); every class must implement \texttt{clone()}
correctly. Watch out for objects owning raw pointers.
\end{complexitybox}

\begin{takeawaybox}
Prototype creates by copying a live object through a virtual
\texttt{clone()}, sidestepping constructors and type switches. Reach for
it when you already have a configured instance and want duplicates, or
when construction is expensive. The one hazard is getting deep-copy
semantics right.
\end{takeawaybox}


% ---- Structural ----
\section{Adapter}
\label{sec:pat-adapter}

\problemstatement{Convert the interface of an existing class into another
interface that clients expect, letting classes that could not otherwise
collaborate -- because of mismatched method names or signatures -- work
together.}

\begin{patternbox}[When You Reach for It]
The trigger is \emph{``I have a class that does the right thing but has
the wrong interface,''} often a third-party or legacy library you cannot
change. When the client expects \texttt{request()} but the library offers
\texttt{specificRequest()}, an Adapter bridges the gap.
\end{patternbox}

\subsection*{Understanding the pattern}

The Adapter implements the interface the client wants and, internally,
holds an instance of the incompatible class (object adapter, via
composition). Each method translates the client's call into one or more
calls on the wrapped object. The client codes against the target
interface; the adapter absorbs the mismatch.

\begin{ideabox}[Wrap the Wrong Interface in the Right One]
Compose the incompatible ``adaptee'' inside a class that presents the
expected interface, forwarding and translating calls. Nothing about the
adaptee or the client changes -- the adapter is the only new, and only
throwaway-if-needed, piece.
\end{ideabox}

\subsection*{C++ implementation}

\begin{lstlisting}
#include <bits/stdc++.h>
using namespace std;

// Target: what the client expects.
struct MediaPlayer {
    virtual void play(const string& file) = 0;
    virtual ~MediaPlayer() = default;
};

// Adaptee: useful, but wrong interface (cannot be modified).
struct LegacyMp4Engine {
    void playMp4(const string& filename) {
        cout << "Playing MP4: " << filename << "\n";
    }
};

// Adapter: presents MediaPlayer, delegates to the legacy engine.
class Mp4Adapter : public MediaPlayer {
    LegacyMp4Engine engine;                    // composed adaptee
public:
    void play(const string& file) override {
        engine.playMp4(file);                  // translate the call
    }
};

int main() {
    unique_ptr<MediaPlayer> player = make_unique<Mp4Adapter>();
    player->play("song.mp4");                  // client uses target API
    return 0;
}
\end{lstlisting}

\begin{complexitybox}[Trade-offs]
\textbf{Pros.} Reuses existing/legacy code without modifying it; isolates
interface-conversion logic in one place; supports Single Responsibility.
\textbf{Cons.} Adds an extra layer; if many methods need adapting the
adapter grows. A sign the underlying interfaces should perhaps be
redesigned if you control both.
\end{complexitybox}

\begin{takeawaybox}
Adapter makes an existing class fit an expected interface by wrapping and
translating -- ``right behaviour, wrong signature.'' Contrast with
Decorator (same interface, added behaviour) and Facade (new simpler
interface over many classes): all three wrap, but for different reasons.
\end{takeawaybox}

\section{Bridge}
\label{sec:pat-bridge}

\problemstatement{Decouple an abstraction from its implementation so the
two can vary independently, avoiding a combinatorial explosion of
subclasses when a type varies along two separate dimensions.}

\begin{patternbox}[When You Reach for It]
The trigger is \emph{``two independent dimensions of variation''} that,
combined by inheritance, would create $m\times n$ subclasses. Shapes
(circle, square) times rendering APIs (OpenGL, DirectX) would need
\texttt{OpenGLCircle}, \texttt{DirectXCircle}, \texttt{OpenGLSquare}\ldots
Bridge turns that product into a sum.
\end{patternbox}

\subsection*{Understanding the pattern}

Split the two dimensions into two class hierarchies and connect them with
a reference (the ``bridge''). The \emph{abstraction} hierarchy (shapes)
holds a pointer to an \emph{implementor} hierarchy (renderers) and
delegates the primitive operations. Now you add shapes and renderers
independently: $m+n$ classes instead of $m\times n$.

\begin{ideabox}[Prefer Composition Across the Two Axes]
When a class varies along two axes, do not multiply them with inheritance.
Give one axis (the abstraction) a member pointing at the other (the
implementation). This is ``favour composition over inheritance'' applied
to prevent a subclass explosion.
\end{ideabox}

\subsection*{C++ implementation}

\begin{lstlisting}
#include <bits/stdc++.h>
using namespace std;

// Implementor axis: how to render.
struct Renderer {
    virtual void renderCircle(float r) = 0;
    virtual ~Renderer() = default;
};
struct OpenGLRenderer : Renderer {
    void renderCircle(float r) override { cout << "OpenGL circle r=" << r << "\n"; }
};
struct DirectXRenderer : Renderer {
    void renderCircle(float r) override { cout << "DirectX circle r=" << r << "\n"; }
};

// Abstraction axis: what to draw. Holds a bridge to the implementor.
class Shape {
protected:
    Renderer& renderer;                        // the bridge
public:
    Shape(Renderer& r) : renderer(r) {}
    virtual void draw() = 0;
    virtual ~Shape() = default;
};
class Circle : public Shape {
    float radius;
public:
    Circle(Renderer& r, float rad) : Shape(r), radius(rad) {}
    void draw() override { renderer.renderCircle(radius); }
};

int main() {
    OpenGLRenderer gl; DirectXRenderer dx;
    Circle c1(gl, 5), c2(dx, 8);               // mix axes freely
    c1.draw();
    c2.draw();
    return 0;
}
\end{lstlisting}

\begin{complexitybox}[Trade-offs]
\textbf{Pros.} Abstraction and implementation evolve independently; avoids
class explosion; swappable implementations at run time. \textbf{Cons.}
More indirection and up-front structure; only pays off when \emph{both}
dimensions genuinely vary. For a single axis it is over-engineering.
\end{complexitybox}

\begin{takeawaybox}
Bridge separates ``what'' from ``how'' into two hierarchies joined by a
reference, turning an $m\times n$ subclass explosion into $m+n$. The
signal is two orthogonal dimensions of change. It looks structurally like
Strategy but its intent is structural decoupling, not swapping an
algorithm.
\end{takeawaybox}

\section{Composite}
\label{sec:pat-composite}

\problemstatement{Compose objects into tree structures to represent
part-whole hierarchies, and let clients treat individual objects and
compositions of objects uniformly through the same interface.}

\begin{patternbox}[When You Reach for It]
The trigger is \emph{``a tree of objects where leaves and containers
should be handled the same way''} -- files and folders, a GUI of widgets
containing widgets, an org chart. When client code should not care whether
it holds one item or a group, use Composite.
\end{patternbox}

\subsection*{Understanding the pattern}

Define one \emph{component} interface implemented by both \emph{leaves}
(no children) and \emph{composites} (hold children, also components). A
composite's operation typically recurses over its children. Because both
share the interface, client code calls, say, \texttt{getSize()} on
anything and the recursion through the tree happens transparently.

\begin{ideabox}[Leaf and Container Share One Interface]
Make the container a component too, holding a list of components. An
operation on a container delegates to its children and aggregates. This
uniformity is the whole point: no \texttt{if (isFolder)} branching in
client code.
\end{ideabox}

\subsection*{C++ implementation}

\begin{lstlisting}
#include <bits/stdc++.h>
using namespace std;

struct FileSystemNode {                        // component
    virtual int size() const = 0;
    virtual void print(string indent = "") const = 0;
    virtual ~FileSystemNode() = default;
};

struct File : FileSystemNode {                 // leaf
    string name; int bytes;
    File(string n, int b) : name(move(n)), bytes(b) {}
    int size() const override { return bytes; }
    void print(string indent) const override {
        cout << indent << name << " (" << bytes << ")\n";
    }
};

struct Folder : FileSystemNode {               // composite
    string name;
    vector<unique_ptr<FileSystemNode>> children;
    Folder(string n) : name(move(n)) {}
    void add(unique_ptr<FileSystemNode> c) { children.push_back(move(c)); }
    int size() const override {
        int total = 0;
        for (auto& c : children) total += c->size();   // recurse
        return total;
    }
    void print(string indent) const override {
        cout << indent << name << "/\n";
        for (auto& c : children) c->print(indent + "  ");
    }
};

int main() {
    auto root = make_unique<Folder>("root");
    root->add(make_unique<File>("a.txt", 100));
    auto sub = make_unique<Folder>("sub");
    sub->add(make_unique<File>("b.txt", 250));
    root->add(move(sub));

    root->print();
    cout << "Total size: " << root->size() << "\n";   // treated uniformly
    return 0;
}
\end{lstlisting}

\begin{complexitybox}[Trade-offs]
\textbf{Pros.} Uniform treatment of leaves and groups; new node types slot
in easily; recursion over the tree is natural. \textbf{Cons.} The shared
interface can become overly general (leaf-only or container-only methods
awkwardly appear on both); type safety of ``only folders have children''
is relaxed.
\end{complexitybox}

\begin{takeawaybox}
Composite models part-whole trees behind a single component interface so
clients treat one and many identically, with container operations
recursing over children. Recognise it wherever a hierarchy of ``things and
groups of things'' must be processed uniformly.
\end{takeawaybox}

\section{Decorator}
\label{sec:pat-decorator}

\problemstatement{Attach additional responsibilities to an object
dynamically, at run time, without subclassing -- providing a flexible
alternative to inheritance for extending behaviour.}

\begin{patternbox}[When You Reach for It]
The trigger is \emph{``add behaviour at run time in combinations,''} where
a subclass per combination would explode (coffee with milk, with sugar,
with milk-and-sugar-and-cream). When features stack in arbitrary
combinations, wrap rather than subclass.
\end{patternbox}

\subsection*{Understanding the pattern}

A decorator implements the \emph{same interface} as the object it wraps
and holds a pointer to a wrapped component of that interface. It forwards
calls to the inner object and adds its own behaviour before or after. Since
a decorator is itself a component, decorators nest: each layer adds one
feature, and you compose features by stacking wrappers at run time.

\begin{ideabox}[Same Interface, Wrap and Add]
Unlike Adapter (changes the interface) a Decorator keeps the interface
identical and augments the behaviour, so wrapped and unwrapped objects are
interchangeable. Stacking decorators composes features additively without
a single new subclass.
\end{ideabox}

\subsection*{C++ implementation}

\begin{lstlisting}
#include <bits/stdc++.h>
using namespace std;

struct Coffee {                                // component
    virtual double cost() const = 0;
    virtual string desc() const = 0;
    virtual ~Coffee() = default;
};

struct Espresso : Coffee {                     // concrete component
    double cost() const override { return 2.0; }
    string desc() const override { return "Espresso"; }
};

struct CoffeeDecorator : Coffee {              // base decorator
    unique_ptr<Coffee> inner;
    CoffeeDecorator(unique_ptr<Coffee> c) : inner(move(c)) {}
};

struct Milk : CoffeeDecorator {
    Milk(unique_ptr<Coffee> c) : CoffeeDecorator(move(c)) {}
    double cost() const override { return inner->cost() + 0.5; }
    string desc() const override { return inner->desc() + " + Milk"; }
};
struct Sugar : CoffeeDecorator {
    Sugar(unique_ptr<Coffee> c) : CoffeeDecorator(move(c)) {}
    double cost() const override { return inner->cost() + 0.2; }
    string desc() const override { return inner->desc() + " + Sugar"; }
};

int main() {
    unique_ptr<Coffee> order =
        make_unique<Sugar>(make_unique<Milk>(make_unique<Espresso>()));
    cout << order->desc() << " = $" << order->cost() << "\n";  // stacked
    return 0;
}
\end{lstlisting}

\begin{complexitybox}[Trade-offs]
\textbf{Pros.} Add/remove responsibilities at run time; avoids a subclass
per feature-combination; each decorator is a small, single-purpose class.
\textbf{Cons.} Many small objects and layers can be hard to debug; order
of wrapping can matter; a deep stack obscures the concrete object.
\end{complexitybox}

\begin{takeawaybox}
Decorator wraps an object in another of the same interface to add
behaviour, and nests to compose features -- the run-time alternative to a
combinatorial subclass tree. Keyword: ``add responsibilities dynamically.''
It is the pattern behind Java's I/O streams and many middleware chains.
\end{takeawaybox}

\section{Facade}
\label{sec:pat-facade}

\problemstatement{Provide a single, simplified interface to a complex
subsystem of many classes, so common tasks become easy while the full
subsystem remains available for advanced use.}

\begin{patternbox}[When You Reach for It]
The trigger is \emph{``clients must call five classes in the right order
just to do one common thing.''} When using a subsystem requires knowing
too much of its internals, a Facade offers a clean front door for the
typical case.
\end{patternbox}

\subsection*{Understanding the pattern}

The Facade is a class that knows the subsystem's classes and orchestrates
them. It exposes a few high-level methods (\texttt{watchMovie()}) that
internally coordinate the amplifier, projector, lights and player. Clients
talk to the Facade; they are shielded from the subsystem's structure and
call order. Crucially the subsystem is not hidden -- power users can still
reach the underlying classes directly.

\begin{ideabox}[A Simple Front Door, Not a Wall]
A Facade reduces coupling between clients and a subsystem by centralising
the common orchestration in one place. It does not add behaviour or
restrict access; it simply makes the easy things easy while leaving the
hard things possible.
\end{ideabox}

\subsection*{C++ implementation}

\begin{lstlisting}
#include <bits/stdc++.h>
using namespace std;

struct Amplifier { void on() { cout << "Amp on\n"; } void setVolume(int v) { cout << "Volume " << v << "\n"; } };
struct Projector { void on() { cout << "Projector on\n"; } void wide() { cout << "Widescreen\n"; } };
struct Lights    { void dim(int p) { cout << "Lights " << p << "%\n"; } };
struct Player    { void play(const string& m) { cout << "Playing " << m << "\n"; } };

class HomeTheaterFacade {                       // the simple front door
    Amplifier amp; Projector projector; Lights lights; Player player;
public:
    void watchMovie(const string& movie) {      // one call orchestrates all
        lights.dim(10);
        projector.on();
        projector.wide();
        amp.on();
        amp.setVolume(7);
        player.play(movie);
    }
};

int main() {
    HomeTheaterFacade theater;
    theater.watchMovie("Inception");            // client does one thing
    return 0;
}
\end{lstlisting}

\begin{complexitybox}[Trade-offs]
\textbf{Pros.} Simplifies client code; decouples clients from subsystem
internals; a natural place to enforce a correct call sequence.
\textbf{Cons.} Can become a ``god object'' if it accumulates too much;
does not prevent misuse of the subsystem directly. Keep it thin -- a
coordinator, not a kitchen sink.
\end{complexitybox}

\begin{takeawaybox}
Facade wraps a complicated subsystem behind a small, task-oriented
interface, lowering the knowledge a client needs. Distinguish from
Adapter (converts one interface to another expected one) -- Facade invents
a \emph{new, simpler} interface over \emph{many} classes.
\end{takeawaybox}

\section{Flyweight}
\label{sec:pat-flyweight}

\problemstatement{Minimise memory use by sharing as much data as possible
between many similar objects, separating the intrinsic (shared,
state-independent) part from the extrinsic (per-instance) part.}

\begin{patternbox}[When You Reach for It]
The trigger is \emph{``millions of near-identical objects blowing up
memory''} -- every character in a document, every tree in a forest, every
bullet in a game. When objects share most of their data, store the shared
part once and pass in what differs.
\end{patternbox}

\subsection*{Understanding the pattern}

Partition an object's state into \textbf{intrinsic} state -- shared,
immutable, independent of context (a glyph's shape and font) -- and
\textbf{extrinsic} state -- unique per use (the glyph's position on the
page). Store one flyweight per distinct intrinsic value in a factory/pool,
hand the same shared flyweight to every user, and pass extrinsic state as
method arguments rather than storing it.

\begin{ideabox}[Share the Intrinsic, Pass the Extrinsic]
Cache and reuse the heavy shared part keyed by its value; supply the light
per-instance part at call time. A thousand `a' characters point to one
shared glyph object; only their coordinates differ and those live in the
document, not in a thousand glyph copies.
\end{ideabox}

\subsection*{C++ implementation}

\begin{lstlisting}
#include <bits/stdc++.h>
using namespace std;

class Glyph {                                   // flyweight: intrinsic state
    char symbol;                                // shared, immutable
public:
    Glyph(char c) : symbol(c) { cout << "Created glyph '" << c << "'\n"; }
    void draw(int x, int y) const {             // extrinsic passed in
        cout << "'" << symbol << "' at (" << x << "," << y << ")\n";
    }
};

class GlyphFactory {                            // ensures sharing
    unordered_map<char, shared_ptr<Glyph>> pool;
public:
    shared_ptr<Glyph> get(char c) {
        auto it = pool.find(c);
        if (it != pool.end()) return it->second;   // reuse
        auto g = make_shared<Glyph>(c);
        pool[c] = g;
        return g;
    }
};

int main() {
    GlyphFactory factory;
    string text = "aab";
    int x = 0;
    for (char c : text) {
        auto g = factory.get(c);                // 'a' created once, shared
        g->draw(x, 0);
        x += 10;
    }
    return 0;
}
\end{lstlisting}

\begin{complexitybox}[Trade-offs]
\textbf{Pros.} Dramatic memory savings when many objects share state;
fewer allocations. \textbf{Cons.} Complexity of splitting state; extrinsic
state must be recomputed or passed each time (a CPU-for-memory trade);
flyweights must be immutable to be safely shared.
\end{complexitybox}

\begin{takeawaybox}
Flyweight shares one immutable intrinsic object among many contexts and
passes the varying extrinsic state as arguments, trading a little CPU for
large memory savings. Recognise it by ``huge number of objects, mostly
identical.'' The pool/factory that guarantees sharing is essential.
\end{takeawaybox}

\section{Proxy}
\label{sec:pat-proxy}

\problemstatement{Provide a surrogate or placeholder for another object to
control access to it -- for lazy initialisation, access control, caching,
logging, or standing in for a remote object.}

\begin{patternbox}[When You Reach for It]
The trigger is \emph{``wrap an object to control how and when it is
used''} without changing its interface: delay an expensive load until
needed (virtual proxy), check permissions (protection proxy), cache
results, or represent something across a network (remote proxy).
\end{patternbox}

\subsection*{Understanding the pattern}

A Proxy implements the \emph{same interface} as the real subject and holds
a reference to it. It intercepts each call, does its extra job (lazy
create, permission check, cache lookup, log), then forwards to the real
object. Because the interface matches, clients cannot tell they are
talking to a proxy rather than the real thing.

\begin{ideabox}[Same Interface, Controlled Access]
Proxy shares Decorator's ``wrap with the same interface'' shape but its
intent is \emph{access control and management}, not adding features. The
canonical use is the virtual proxy: construct the heavy real object only
on first genuine use.
\end{ideabox}

\subsection*{C++ implementation}

\begin{lstlisting}
#include <bits/stdc++.h>
using namespace std;

struct Image {
    virtual void display() = 0;
    virtual ~Image() = default;
};

class RealImage : public Image {                // expensive real subject
    string filename;
public:
    RealImage(const string& f) : filename(f) {
        cout << "Loading " << filename << " from disk\n";   // heavy
    }
    void display() override { cout << "Displaying " << filename << "\n"; }
};

class ImageProxy : public Image {               // virtual proxy
    string filename;
    unique_ptr<RealImage> real;                 // created lazily
public:
    ImageProxy(const string& f) : filename(f) {}
    void display() override {
        if (!real) real = make_unique<RealImage>(filename);  // load on demand
        real->display();
    }
};

int main() {
    ImageProxy img("photo.png");                // no disk load yet
    cout << "Proxy created, nothing loaded.\n";
    img.display();                              // loads now, then displays
    img.display();                              // reuses, no reload
    return 0;
}
\end{lstlisting}

\begin{complexitybox}[Trade-offs]
\textbf{Pros.} Adds access control, laziness, caching or remoting
transparently; the client interface is unchanged. \textbf{Cons.} Extra
indirection and latency on first access; a poorly designed proxy can hide
performance costs or introduce inconsistency between proxy and subject.
\end{complexitybox}

\begin{takeawaybox}
Proxy stands in for a real object behind an identical interface to control
access -- lazy loading, protection, caching, or remoting. It looks like
Decorator but the intent differs: Decorator \emph{adds behaviour}, Proxy
\emph{governs access}. Name the proxy variant (virtual, protection,
remote) when you use it.
\end{takeawaybox}


% ---- Behavioural ----
\section{Chain of Responsibility}
\label{sec:pat-chain}

\problemstatement{Pass a request along a chain of handlers; each handler
decides either to process the request or to forward it to the next,
decoupling the sender from the specific receiver that ends up handling it.}

\begin{patternbox}[When You Reach for It]
The trigger is \emph{``a request that several handlers might process, and
we don't know in advance which one will.''} Approval workflows, event
bubbling, logging levels, and middleware pipelines all fit: try each
handler in turn until one takes it.
\end{patternbox}

\subsection*{Understanding the pattern}

Each handler holds a pointer to the next and implements a
\texttt{handle(request)}: if it can deal with the request it does;
otherwise it delegates to the next in line. The client fires the request
at the head of the chain and never needs to know which link resolves it.
Handlers can be reordered or inserted without touching the sender.

\begin{ideabox}[Handle or Pass Along]
Give every handler the same interface and a ``next'' link. The decision
``mine or not mine'' is local to each handler; the chain's shape is
configured externally. This decouples \emph{who sends} from \emph{who
handles}.
\end{ideabox}

\subsection*{C++ implementation}

\begin{lstlisting}
#include <bits/stdc++.h>
using namespace std;

class Approver {
protected:
    Approver* next = nullptr;
public:
    Approver* setNext(Approver* n) { next = n; return n; }
    virtual void handle(int amount) {
        if (next) next->handle(amount);        // default: pass along
        else cout << "No one could approve " << amount << "\n";
    }
    virtual ~Approver() = default;
};

class Manager : public Approver {
public:
    void handle(int amount) override {
        if (amount <= 1000) cout << "Manager approves " << amount << "\n";
        else Approver::handle(amount);
    }
};
class Director : public Approver {
public:
    void handle(int amount) override {
        if (amount <= 10000) cout << "Director approves " << amount << "\n";
        else Approver::handle(amount);
    }
};

int main() {
    Manager mgr; Director dir;
    mgr.setNext(&dir);                          // build the chain
    mgr.handle(500);                            // Manager
    mgr.handle(5000);                           // passed to Director
    mgr.handle(50000);                          // unhandled
    return 0;
}
\end{lstlisting}

\begin{complexitybox}[Trade-offs]
\textbf{Pros.} Decouples sender and receiver; handlers are added,
removed, or reordered freely; each handler is simple. \textbf{Cons.} A
request may fall off the end unhandled; harder to trace which handler
acted; a long chain adds latency. Ensure a terminal/default handler.
\end{complexitybox}

\begin{takeawaybox}
Chain of Responsibility lets a request travel handler-to-handler until one
processes it, decoupling the caller from the resolver. Recognise it in
approval hierarchies and middleware. The key risk to mention: unhandled
requests, so design an end-of-chain default.
\end{takeawaybox}

\section{Command}
\label{sec:pat-command}

\problemstatement{Encapsulate a request as an object, so requests can be
parameterised, queued, logged, and undone -- decoupling the object that
invokes an operation from the one that knows how to perform it.}

\begin{patternbox}[When You Reach for It]
The triggers are \emph{``undo/redo,''} \emph{``queue or schedule
operations,''} and \emph{``parameterise a button/menu with an action.''}
When an action must be treated as a first-class value -- stored, passed,
reversed -- wrap it in a Command object.
\end{patternbox}

\subsection*{Understanding the pattern}

A Command object exposes \texttt{execute()} (and often \texttt{undo()})
and holds everything the action needs: a reference to the receiver and any
arguments. An invoker (a button, a queue, a macro) triggers commands
without knowing what they do. Because each command is an object, you can
keep a history stack for undo, replay a log, or batch commands into a
macro.

\begin{ideabox}[Turn a Method Call into an Object]
Reifying ``call this method with these arguments'' as an object is what
enables undo, queuing, logging and macros -- none of which are possible
with a bare function call. The invoker depends only on the Command
interface.
\end{ideabox}

\subsection*{C++ implementation}

\begin{lstlisting}
#include <bits/stdc++.h>
using namespace std;

struct Light {                                  // receiver
    bool on = false;
    void turnOn()  { on = true;  cout << "Light on\n"; }
    void turnOff() { on = false; cout << "Light off\n"; }
};

struct Command {
    virtual void execute() = 0;
    virtual void undo() = 0;
    virtual ~Command() = default;
};

struct LightOnCommand : Command {
    Light& light;
    LightOnCommand(Light& l) : light(l) {}
    void execute() override { light.turnOn(); }
    void undo() override { light.turnOff(); }
};

class RemoteControl {                            // invoker with undo
    vector<unique_ptr<Command>> history;
public:
    void submit(unique_ptr<Command> cmd) {
        cmd->execute();
        history.push_back(move(cmd));
    }
    void undoLast() {
        if (!history.empty()) {
            history.back()->undo();
            history.pop_back();
        }
    }
};

int main() {
    Light light;
    RemoteControl remote;
    remote.submit(make_unique<LightOnCommand>(light));
    remote.undoLast();                          // reverses it
    return 0;
}
\end{lstlisting}

\begin{complexitybox}[Trade-offs]
\textbf{Pros.} Enables undo/redo, queuing, logging, and macros; decouples
invoker from receiver; commands are composable. \textbf{Cons.} A class per
action can proliferate; undo requires each command to capture enough state
to reverse itself, which can be non-trivial.
\end{complexitybox}

\begin{takeawaybox}
Command packages an action plus its receiver and arguments into an object
with \texttt{execute}/\texttt{undo}, unlocking undo stacks, queues, and
macros. The interview keyword is almost always ``undo/redo.'' It is the
backbone of editors and task schedulers.
\end{takeawaybox}

\section{Iterator}
\label{sec:pat-iterator}

\problemstatement{Provide a way to access the elements of an aggregate
object sequentially without exposing its underlying representation, so
clients traverse a collection the same way regardless of how it is stored.}

\begin{patternbox}[When You Reach for It]
The trigger is \emph{``traverse a collection without revealing whether
it's an array, list, or tree.''} When you want uniform iteration and the
freedom to change the container's internals, expose an iterator instead of
raw structure.
\end{patternbox}

\subsection*{Understanding the pattern}

Define an iterator interface with \texttt{hasNext()} and \texttt{next()}.
The collection produces an iterator that knows how to walk \emph{its}
particular structure. Client loops depend only on the iterator interface,
so swapping a vector for a linked list -- or offering multiple traversal
orders -- does not touch client code. This is the pattern the C++ STL and
range-based \texttt{for} are built on.

\begin{ideabox}[Externalise Traversal State]
Move the ``where am I in the collection'' cursor out of the client and
into a dedicated iterator object. Multiple independent iterators can then
traverse the same collection at once, and the collection's storage stays
hidden.
\end{ideabox}

\subsection*{C++ implementation}

\begin{lstlisting}
#include <bits/stdc++.h>
using namespace std;

template <typename T>
struct Iterator {
    virtual bool hasNext() const = 0;
    virtual T next() = 0;
    virtual ~Iterator() = default;
};

class IntCollection {                           // hides its storage
    vector<int> data;
public:
    void add(int x) { data.push_back(x); }

    class Iter : public Iterator<int> {
        const vector<int>& ref; size_t pos = 0;
    public:
        Iter(const vector<int>& r) : ref(r) {}
        bool hasNext() const override { return pos < ref.size(); }
        int next() override { return ref[pos++]; }
    };

    unique_ptr<Iterator<int>> iterator() const {
        return make_unique<Iter>(data);
    }
};

int main() {
    IntCollection c;
    c.add(10); c.add(20); c.add(30);
    auto it = c.iterator();
    while (it->hasNext()) cout << it->next() << " ";   // storage hidden
    cout << "\n";
    return 0;
}
\end{lstlisting}

\begin{complexitybox}[Trade-offs]
\textbf{Pros.} Uniform traversal; hides representation; supports multiple
simultaneous and multiple-order traversals; simplifies the collection's
own interface. \textbf{Cons.} Extra objects for simple cases; an iterator
can be invalidated if the collection is modified during traversal.
\end{complexitybox}

\begin{takeawaybox}
Iterator externalises traversal into an object with
\texttt{hasNext}/\texttt{next}, decoupling clients from a collection's
internal structure. It is so fundamental that most languages bake it in
(STL iterators, \texttt{for-each}). Mention iterator invalidation as the
classic pitfall.
\end{takeawaybox}

\section{Mediator}
\label{sec:pat-mediator}

\problemstatement{Define an object that encapsulates how a set of objects
interact, so they refer to the mediator instead of to each other -- turning
a tangle of many-to-many references into a hub-and-spoke.}

\begin{patternbox}[When You Reach for It]
The trigger is \emph{``every component talks to every other component''} --
an $n^2$ web of direct references (a dialog where each widget reacts to
others, a chat room, air-traffic control). When coupling between peers
grows quadratic, route it through a mediator.
\end{patternbox}

\subsection*{Understanding the pattern}

Instead of colleagues holding references to one another, each holds a
reference to a single \emph{mediator}. When a colleague changes, it
notifies the mediator, which decides who else needs to know and calls
them. The colleagues become simple and reusable; all the interaction logic
concentrates in one place -- the mediator -- which is the only object that
knows the participants.

\begin{ideabox}[Centralise the Many-to-Many]
Replace peer-to-peer links with a hub. Colleagues know only the mediator;
the mediator knows the colleagues. This converts $O(n^2)$ coupling into
$O(n)$ and makes interaction rules explicit and changeable in one spot.
\end{ideabox}

\subsection*{C++ implementation}

\begin{lstlisting}
#include <bits/stdc++.h>
using namespace std;

class User;
struct ChatRoom {                               // mediator
    virtual void send(const string& msg, User* from) = 0;
    virtual ~ChatRoom() = default;
};

class User {
    string name;
    ChatRoom* room;                             // knows only the mediator
public:
    User(string n, ChatRoom* r) : name(move(n)), room(r) {}
    const string& getName() const { return name; }
    void send(const string& msg);
    void receive(const string& msg, const string& from) {
        cout << name << " got from " << from << ": " << msg << "\n";
    }
};

class Room : public ChatRoom {
    vector<User*> users;
public:
    void join(User* u) { users.push_back(u); }
    void send(const string& msg, User* from) override {
        for (User* u : users)
            if (u != from) u->receive(msg, from->getName());   // hub routes
    }
};

void User::send(const string& msg) { room->send(msg, this); }

int main() {
    Room room;
    User a("Alice", &room), b("Bob", &room), c("Carol", &room);
    room.join(&a); room.join(&b); room.join(&c);
    a.send("hi all");                           // no direct peer links
    return 0;
}
\end{lstlisting}

\begin{complexitybox}[Trade-offs]
\textbf{Pros.} Removes tight peer-to-peer coupling; interaction logic
lives in one place; colleagues become reusable. \textbf{Cons.} The
mediator can swell into a complex ``god object'' as rules accumulate --
you have centralised complexity, not removed it.
\end{complexitybox}

\begin{takeawaybox}
Mediator replaces a many-to-many mesh of collaborators with a hub that
owns the interaction rules, cutting coupling from $O(n^2)$ to $O(n)$.
Watch the mediator itself for bloat. Classic examples: chat rooms, UI
dialog coordination, air-traffic control.
\end{takeawaybox}

\section{Memento}
\label{sec:pat-memento}

\problemstatement{Capture and externalise an object's internal state
without violating encapsulation, so the object can be restored to that
state later -- the mechanism behind save points and undo.}

\begin{patternbox}[When You Reach for It]
The trigger is \emph{``save and restore an object's state''} -- snapshots,
checkpoints, undo history. When you must be able to roll an object back but
do not want to expose its internals to the code that stores the snapshot,
use a Memento.
\end{patternbox}

\subsection*{Understanding the pattern}

Three roles: the \textbf{Originator} (whose state we snapshot) creates a
\textbf{Memento} holding a copy of its state and can later restore from
one; the \textbf{Caretaker} keeps mementos (e.g.\ on an undo stack) but
treats them as opaque -- it never reads their contents. Encapsulation is
preserved because only the Originator can read or write a memento's
internals.

\begin{ideabox}[Opaque Snapshots the Caretaker Cannot Read]
The Originator packages its own state into a memento and unpacks it on
restore; the Caretaker only stores and hands mementos back. This division
gives you undo without leaking private state to the code that manages
history.
\end{ideabox}

\subsection*{C++ implementation}

\begin{lstlisting}
#include <bits/stdc++.h>
using namespace std;

class Memento {                                 // opaque snapshot
    friend class Editor;                        // only Editor sees inside
    string state;
    Memento(string s) : state(move(s)) {}
};

class Editor {                                  // originator
    string content;
public:
    void type(const string& text) { content += text; }
    string show() const { return content; }
    Memento save() const { return Memento(content); }      // snapshot
    void restore(const Memento& m) { content = m.state; }  // roll back
};

int main() {
    Editor editor;
    editor.type("Hello");
    Memento checkpoint = editor.save();         // caretaker holds this
    editor.type(", World");
    cout << "Before undo: " << editor.show() << "\n";
    editor.restore(checkpoint);                 // back to "Hello"
    cout << "After undo:  " << editor.show() << "\n";
    return 0;
}
\end{lstlisting}

\begin{complexitybox}[Trade-offs]
\textbf{Pros.} Clean save/restore without breaking encapsulation; pairs
naturally with Command for undo. \textbf{Cons.} Snapshots can be
memory-heavy if state is large or saved often; deciding what to include in
a memento and how frequently to snapshot needs care (consider
incremental/diff mementos).
\end{complexitybox}

\begin{takeawaybox}
Memento snapshots an object's state into an opaque object that only the
originator can interpret, enabling undo and checkpoints while keeping
internals private. It frequently teams with Command: the command stores a
memento to reverse itself. Cost to flag: memory for many/large snapshots.
\end{takeawaybox}

\section{Observer}
\label{sec:pat-observer}

\problemstatement{Define a one-to-many dependency so that when one object
(the subject) changes state, all its dependents (observers) are notified
and updated automatically -- the foundation of event systems and
publish/subscribe.}

\begin{patternbox}[When You Reach for It]
The trigger is \emph{``when X changes, several other things must react''} --
a spreadsheet cell and its charts, a stock price and its displays, a model
and its views. When one source has many interested listeners that come and
go, use Observer.
\end{patternbox}

\subsection*{Understanding the pattern}

The \textbf{Subject} maintains a list of \textbf{Observers} and offers
\texttt{subscribe}/\texttt{unsubscribe}. On a relevant change it calls
\texttt{notify()}, which invokes \texttt{update()} on every observer. The
subject knows only the observer \emph{interface}, so observers can be added
or removed at run time without the subject knowing their concrete types.
This is the classic Model-View decoupling.

\begin{ideabox}[Publish to a List of Subscribers]
Keep subscribers behind an interface and iterate them on change. Loose
coupling is the point: the subject broadcasts an event; who listens, and
how many, is not its concern and can vary dynamically.
\end{ideabox}

\subsection*{C++ implementation}

\begin{lstlisting}
#include <bits/stdc++.h>
using namespace std;

struct Observer {
    virtual void update(int price) = 0;
    virtual ~Observer() = default;
};

class Stock {                                   // subject
    vector<Observer*> observers;
    int price = 0;
public:
    void subscribe(Observer* o)   { observers.push_back(o); }
    void unsubscribe(Observer* o) {
        observers.erase(remove(observers.begin(), observers.end(), o),
                        observers.end());
    }
    void setPrice(int p) {
        price = p;
        for (Observer* o : observers) o->update(price);   // notify all
    }
};

struct Display : Observer {
    string name;
    Display(string n) : name(move(n)) {}
    void update(int price) override {
        cout << name << " sees price " << price << "\n";
    }
};

int main() {
    Stock stock;
    Display a("Phone"), b("Web");
    stock.subscribe(&a); stock.subscribe(&b);
    stock.setPrice(100);                        // both notified
    stock.unsubscribe(&a);
    stock.setPrice(120);                        // only Web notified
    return 0;
}
\end{lstlisting}

\begin{complexitybox}[Trade-offs]
\textbf{Pros.} Loose coupling between subject and observers; observers
added/removed at run time; supports broadcast/event architectures.
\textbf{Cons.} Notification order is unspecified; risk of update storms or
cycles; a dangling observer that forgot to unsubscribe causes bugs or
leaks (use weak references or careful lifetime management).
\end{complexitybox}

\begin{takeawaybox}
Observer wires a one-to-many ``notify on change'' relationship behind an
observer interface, the core of MVC and pub/sub. Interview keyword: ``when
this changes, update all those.'' The pitfalls to name are lapsed-listener
leaks and unspecified notification order.
\end{takeawaybox}

\section{State}
\label{sec:pat-state}

\problemstatement{Allow an object to alter its behaviour when its internal
state changes, so it appears to change its class -- replacing sprawling
conditionals on a state variable with a set of state objects.}

\begin{patternbox}[When You Reach for It]
The trigger is \emph{``behaviour depends on a mode/state, and the same
method does different things in each''} -- a vending machine, a media
player (playing/paused/stopped), a TCP connection, an order lifecycle.
When methods are riddled with \texttt{if (state == ...)}, extract states
into classes.
\end{patternbox}

\subsection*{Understanding the pattern}

Each state becomes a class implementing a common \texttt{State} interface;
the context holds a pointer to the current state and delegates behaviour to
it. A state's methods perform the state-specific action and can transition
the context to another state. Adding a new state is a new class, not
another branch in every method -- Open/Closed for state machines.

\begin{ideabox}[One Class per State, Delegate and Transition]
Turn each mode into an object that knows how to behave \emph{and} which
state comes next. The context just forwards calls to its current state
object. The scattered \texttt{switch(state)} logic collapses into
localised, self-contained state classes.
\end{ideabox}

\subsection*{C++ implementation}

\begin{lstlisting}
#include <bits/stdc++.h>
using namespace std;

class Fan;
struct State {
    virtual void pull(Fan& fan) = 0;            // the chain-pull action
    virtual string name() const = 0;
    virtual ~State() = default;
};

class Fan {
    unique_ptr<State> state;
public:
    Fan(unique_ptr<State> s) : state(move(s)) {}
    void setState(unique_ptr<State> s) { state = move(s); }
    void pull() { state->pull(*this); }
    string status() const { return state->name(); }
};

struct Off; struct Low;
struct Low : State {
    void pull(Fan& fan) override;
    string name() const override { return "Low"; }
};
struct Off : State {
    void pull(Fan& fan) override { fan.setState(make_unique<Low>()); }
    string name() const override { return "Off"; }
};
void Low::pull(Fan& fan) { fan.setState(make_unique<Off>()); }

int main() {
    Fan fan(make_unique<Off>());
    cout << fan.status() << "\n";               // Off
    fan.pull(); cout << fan.status() << "\n";   // Low
    fan.pull(); cout << fan.status() << "\n";   // Off
    return 0;
}
\end{lstlisting}

\begin{complexitybox}[Trade-offs]
\textbf{Pros.} Eliminates large state conditionals; each state is
isolated and testable; transitions are explicit; new states are added
without editing others. \textbf{Cons.} More classes; transition logic is
spread across state classes, so the overall state machine can be harder to
see at a glance (a transition table can complement it).
\end{complexitybox}

\begin{takeawaybox}
State replaces ``behaviour depends on a mode'' conditionals with one class
per state that both acts and decides the next state. Structurally it is
Strategy, but the intent differs: State objects \emph{swap themselves} as
the machine transitions, whereas a Strategy is chosen by the client.
\end{takeawaybox}

\section{Strategy}
\label{sec:pat-strategy}

\problemstatement{Define a family of interchangeable algorithms,
encapsulate each one, and make them swappable at run time so the algorithm
varies independently of the client that uses it.}

\begin{patternbox}[When You Reach for It]
The trigger is \emph{``several ways to do the same thing, chosen at run
time''} -- different sort orders, payment methods, compression schemes,
route-finding algorithms. When you catch yourself writing \texttt{if
(mode == A) doA(); else doB();}, extract the modes into strategy objects.
\end{patternbox}

\subsection*{Understanding the pattern}

Encapsulate each algorithm behind a common interface. The context holds a
pointer to a strategy and delegates the varying step to it, without knowing
which concrete algorithm it is. Clients pick and inject a strategy; they
can change it at run time. This is the Open/Closed Principle for
behaviour: add a new algorithm as a new class, edit nothing existing.

\begin{ideabox}[Inject the Algorithm as an Object]
Pull the part that varies into a strategy object and hand it to the
context. The context is written once against the strategy interface; the
concrete behaviour is a plug-in chosen by the client. This is composition
replacing a conditional.
\end{ideabox}

\subsection*{C++ implementation}

\begin{lstlisting}
#include <bits/stdc++.h>
using namespace std;

struct PaymentStrategy {
    virtual void pay(int amount) = 0;
    virtual ~PaymentStrategy() = default;
};
struct CreditCard : PaymentStrategy {
    void pay(int amt) override { cout << "Paid " << amt << " by card\n"; }
};
struct UpiPayment : PaymentStrategy {
    void pay(int amt) override { cout << "Paid " << amt << " via UPI\n"; }
};

class Checkout {                                // context
    unique_ptr<PaymentStrategy> strategy;
public:
    void setStrategy(unique_ptr<PaymentStrategy> s) { strategy = move(s); }
    void pay(int amount) { strategy->pay(amount); }   // delegate varying step
};

int main() {
    Checkout checkout;
    checkout.setStrategy(make_unique<CreditCard>());
    checkout.pay(500);
    checkout.setStrategy(make_unique<UpiPayment>());  // swap at run time
    checkout.pay(300);
    return 0;
}
\end{lstlisting}

\begin{complexitybox}[Trade-offs]
\textbf{Pros.} Algorithms are interchangeable at run time; removes
conditionals; each algorithm is isolated and testable; supports
Open/Closed. \textbf{Cons.} Clients must know the strategies exist to
choose one; more objects/classes; overkill when the behaviour never varies.
\end{complexitybox}

\begin{takeawaybox}
Strategy encapsulates interchangeable algorithms behind one interface and
lets the client inject the choice, replacing behavioural conditionals with
composition. It is the single most common behavioural pattern in
interviews. Contrast with State: same shape, but the \emph{client} picks a
Strategy while a State machine swaps its own state.
\end{takeawaybox}

\section{Template Method}
\label{sec:pat-template-method}

\problemstatement{Define the skeleton of an algorithm in a base-class
method, deferring some steps to subclasses, so subclasses can redefine
certain steps without changing the algorithm's overall structure.}

\begin{patternbox}[When You Reach for It]
The trigger is \emph{``the same overall procedure with a few steps that
differ''} -- data importers that all parse/validate/save but read
different formats, report generators, test set-up/tear-down. When the
outline is fixed but a few steps vary, fix the outline and vary the steps.
\end{patternbox}

\subsection*{Understanding the pattern}

A base class provides a non-virtual \texttt{templateMethod()} that calls a
sequence of steps in a fixed order. Some steps are concrete (shared), some
are \texttt{virtual}/abstract ``hooks'' that subclasses override. The
invariant sequence lives in one place; only the variable steps are
subclassed. This inverts control: the base class calls down into the
subclass (``don't call us, we'll call you'').

\begin{ideabox}[Fix the Skeleton, Override the Steps]
Keep the algorithm's order in a base method and expose just the varying
steps as overridable hooks. Subclasses cannot accidentally reorder or skip
the algorithm -- they only fill in the blanks.
\end{ideabox}

\subsection*{C++ implementation}

\begin{lstlisting}
#include <bits/stdc++.h>
using namespace std;

class DataProcessor {
public:
    void process() {                            // template method (fixed order)
        read();
        auto data = parse();
        cout << "Processing " << data << " rows\n";
        save();
    }
    virtual ~DataProcessor() = default;
protected:
    virtual void read() = 0;                     // varying steps
    virtual int  parse() = 0;
    void save() { cout << "Saved to DB\n"; }     // shared step
};

class CsvProcessor : public DataProcessor {
protected:
    void read() override { cout << "Read CSV file\n"; }
    int  parse() override { return 42; }
};

int main() {
    CsvProcessor p;
    p.process();                                // skeleton runs, steps vary
    return 0;
}
\end{lstlisting}

\begin{complexitybox}[Trade-offs]
\textbf{Pros.} Removes duplicated algorithm structure; enforces a fixed
sequence; subclasses supply only what differs. \textbf{Cons.} Relies on
inheritance (less flexible than composition); the base--subclass contract
about which hooks to override can be subtle; changing the skeleton affects
all subclasses.
\end{complexitybox}

\begin{takeawaybox}
Template Method locks an algorithm's structure in a base method and lets
subclasses override only the variable steps -- inheritance-based inversion
of control. Its composition-based cousin is Strategy: Template Method
varies steps via subclassing, Strategy varies the whole algorithm via an
injected object.
\end{takeawaybox}

\section{Visitor}
\label{sec:pat-visitor}

\problemstatement{Represent an operation to be performed on the elements of
an object structure, letting you define a new operation without changing
the classes of the elements it operates on.}

\begin{patternbox}[When You Reach for It]
The trigger is \emph{``add new operations over a stable set of element
types''} -- an AST you want to pretty-print, type-check, and compile,
without editing every node class each time. When element types rarely
change but operations keep being added, Visitor centralises each operation.
\end{patternbox}

\subsection*{Understanding the pattern}

Each element exposes \texttt{accept(visitor)}, which calls back
\texttt{visitor.visitConcreteType(*this)} -- a technique called
\textbf{double dispatch}, because the executed method depends on both the
element's type and the visitor's type. A visitor bundles one operation
across all element types as a set of \texttt{visit} overloads. Adding an
operation is a new visitor class; the element classes are untouched.

\begin{ideabox}[Double Dispatch Moves Operations Out of the Elements]
\texttt{accept} resolves the element's concrete type; the matching
\texttt{visit} overload resolves the operation. Together they route to the
right (element, operation) pair without a type switch, so operations live
in visitors rather than being smeared across every element class.
\end{ideabox}

\subsection*{C++ implementation}

\begin{lstlisting}
#include <bits/stdc++.h>
using namespace std;

struct Circle; struct Square;
struct Visitor {
    virtual void visit(Circle&) = 0;
    virtual void visit(Square&) = 0;
    virtual ~Visitor() = default;
};

struct Shape { virtual void accept(Visitor& v) = 0; virtual ~Shape() = default; };
struct Circle : Shape {
    double r = 2;
    void accept(Visitor& v) override { v.visit(*this); }   // double dispatch
};
struct Square : Shape {
    double side = 3;
    void accept(Visitor& v) override { v.visit(*this); }
};

struct AreaVisitor : Visitor {                  // a new operation, no shape edits
    void visit(Circle& c) override { cout << "Circle area " << 3.14159 * c.r * c.r << "\n"; }
    void visit(Square& s) override { cout << "Square area " << s.side * s.side << "\n"; }
};

int main() {
    vector<unique_ptr<Shape>> shapes;
    shapes.push_back(make_unique<Circle>());
    shapes.push_back(make_unique<Square>());
    AreaVisitor area;
    for (auto& s : shapes) s->accept(area);     // right visit() per type
    return 0;
}
\end{lstlisting}

\begin{complexitybox}[Trade-offs]
\textbf{Pros.} Add new operations easily (one new visitor); gathers a
scattered operation into one class; keeps element classes clean.
\textbf{Cons.} Adding a new \emph{element type} forces editing every
visitor -- the exact opposite trade-off to most patterns. Use only when
element types are stable and operations grow.
\end{complexitybox}

\begin{takeawaybox}
Visitor externalises operations over an object structure via double
dispatch, so new operations are new visitor classes with the elements
untouched -- at the cost of making new element types expensive. It is the
inverse trade-off of a type switch, and the go-to for AST/compiler passes.
\end{takeawaybox}

